\section{Quantitative Results}

This section evaluates the model quantitatively and connects each theoretical mechanism to simulated outcomes.

\subsection{Calibration and Setup}

The model is simulated across a grid of overvaluation levels. Parameters are chosen to ensure three properties: (i) foreign-currency demand increases monotonically in overvaluation, (ii) both crisis and non-crisis outcomes occur with positive probability, and (iii) stablecoins affect both access costs and information precision.

\subsection{Foreign-Currency Demand}

Foreign-currency demand increases monotonically with overvaluation. The increase is driven initially by hedge demand and later by run demand as crisis beliefs rise. Stablecoin usage rises with overvaluation and accounts for a growing share of total FX demand.

\subsection{Crisis Probability}

Crisis probability increases with overvaluation in both the cash-only and stablecoin economies. However, the presence of stablecoins shifts the crisis probability upward at all levels of overvaluation. This reflects both increased demand and stronger coordination.

\subsection{Welfare}

Welfare exhibits a non-monotonic pattern. In low and moderate overvaluation states, stablecoins improve welfare by lowering access costs and increasing consumption smoothing. In high-overvaluation states, stablecoins reduce welfare by increasing crisis probability and amplifying losses in crisis states.

\subsection{Decomposition}

Counterfactual simulations shutting down individual channels show that both the cost-reduction channel and the information channel are necessary. Removing the cost channel eliminates welfare gains, while removing the information channel eliminates crisis amplification.

These results confirm that the model’s mechanisms operate as predicted and jointly generate the patterns observed in the data.
